\section{Introduction}

Knowledge-intensive language-model systems can answer from parameters, retrieve external evidence, cache reusable facts, or update themselves temporarily or durably. Prior work usually studies retrieval selection \citep{asai2024selfrag,yan2024crag} or editing mechanisms \citep{wang2024wise,yu2023melo,zhang2024keStudy} in isolation. We study the controller problem instead: when should one policy choose among these operations under one objective?

Our router operates over six actions: \texttt{PARAM\_ONLY}, \texttt{READ\_MEMORY}, \texttt{RETRIEVE}, \texttt{WRITE\_MEMORY}, \texttt{FAST\_ADAPT}, and \texttt{CONSOLIDATE}. It optimizes
\[
r_t = q_t - c_t - f_t + \beta V_{\text{future}}(x_t, a_t),
\]
where the future-value term rewards actions whose utility is expected to recur later.

\section{Benchmarks}

We use the same lightweight backbone, \texttt{google/flan-t5-small} \citep{longpre2023flan}, across both frozen experiments. PopQA \citep{mallen2022trust} is the recurring-memory benchmark, evaluated on a cache-first recurring-heavy 500-example slice. The freshness benchmark is a deterministic bundled \texttt{FreshQA}-style snapshot inspired by FreshLLMs/FreshQA \citep{vu2024freshllms}; it is included with the submission materials and is \emph{not} presented as a public benchmark run.

\section{Main Results}

\begin{table}[h]
\centering
\small
\begin{tabular}{lrrrrrr}
\toprule
Mode & Quality & Latency & Retrieve & Read & Write & Adapt \\
\midrule
param\_only & 0.064 & 0.033 & 0 & 0 & 0 & 0 \\
always\_retrieve & 0.178 & 0.360 & 500 & 0 & 0 & 0 \\
retrieve\_gate & 0.178 & 0.364 & 500 & 0 & 0 & 0 \\
router & 0.178 & 0.334 & 267 & 327 & 156 & 0 \\
\bottomrule
\end{tabular}


\caption{PopQA main result from the frozen post-fix bundle.}
\label{tab:w-popqa}
\end{table}

\begin{table}[h]
\centering
\small
\begin{tabular}{lrrrrrr}
\toprule
Mode & Overall & Recurring & Non-rec. & Rec. Retrieve & Rec. Read & Rec. Write \\
\midrule
param\_only & 0.064 & 0.056 & 0.093 & 0.000 & 0.000 & 0.000 \\
always\_retrieve & 0.178 & 0.181 & 0.167 & 1.000 & 0.000 & 0.000 \\
retrieve\_gate & 0.178 & 0.181 & 0.167 & 1.000 & 0.000 & 0.000 \\
router & 0.178 & 0.181 & 0.167 & 0.406 & 0.834 & 0.398 \\
\bottomrule
\end{tabular}


\caption{PopQA recurring versus non-recurring breakdown.}
\label{tab:w-popqa-split}
\end{table}

On PopQA, the controller now matches retrieval-level quality overall while preserving recurring-case quality and reducing recurring retrieval from `1.000` to `0.406`. This is the recurring-memory claim: once a fact is reusable, routing can replace repeated retrieval with controlled memory use without sacrificing quality.

\begin{table}[h]
\centering
\tiny
\begin{tabular}{lrrrrrrrrrr}
\toprule
Mode & Quality & Stale & Latency & Retrieve & Read & Write & Adapt & Cons. & Rollback & Forget \\
\midrule
param\_only & 0.133 & 0.867 & 0.200 & 0 & 0 & 0 & 0 & 0 & 0 & 0.000 \\
always\_retrieve & 0.133 & 0.867 & 0.564 & 13 & 0 & 0 & 0 & 0 & 0 & 0.000 \\
memory\_only & 0.600 & 0.400 & 0.357 & 3 & 8 & 3 & 0 & 0 & 0 & 0.000 \\
fast\_adapt\_only & 0.800 & 0.200 & 0.536 & 7 & 0 & 0 & 7 & 0 & 0 & 0.000 \\
router & 0.867 & 0.133 & 0.591 & 8 & 6 & 3 & 4 & 2 & 1 & 0.000 \\
\bottomrule
\end{tabular}


\caption{Freshness main result from the frozen \texttt{freshness\_v1} bundle.}
\label{tab:w-freshness}
\end{table}

On freshness, the full router reaches `0.867` answer quality and `0.133` stale-answer rate, outperforming retrieval-only and memory-only while using \texttt{FAST\_ADAPT} meaningfully, keeping \texttt{CONSOLIDATE} rare, and exercising rollback. This is the updated-knowledge claim: the controller can decide when updated knowledge should be retrieved, written to memory, adapted temporarily, or promoted durably.

\begin{table}[h]
\centering
\small
\begin{tabular}{lrrrr}
\toprule
Benchmark / metric & Full & Action ablation 1 & Action ablation 2 & No $V_{\text{future}}$ \\
\midrule
PopQA overall quality & 0.178 & 0.132 (no read) & 0.078 (no write) & 0.082 \\
Freshness answer quality & 0.867 & 0.667 (no adapt) & 0.733 (no cons.) & 0.733 \\
Freshness stale-answer rate & 0.133 & 0.333 (no adapt) & 0.267 (no cons.) & 0.267 \\
\bottomrule
\end{tabular}


\caption{Compact ablation summary across the two frozen benchmarks.}
\label{tab:w-ablations}
\end{table}

The ablations show that the gains are not incidental. On PopQA, removing \texttt{READ\_MEMORY}, \texttt{WRITE\_MEMORY}, or $V_{\text{future}}$ degrades the recurring-memory result. On freshness, removing \texttt{FAST\_ADAPT}, \texttt{CONSOLIDATE}, or $V_{\text{future}}$ degrades updated-knowledge behavior.

\section{Error Analysis And Transparency}

The frozen PopQA audit shows that the main pre-fix weakness was avoidable non-recurring routing error; the final post-fix bundle removes that harm and leaves mostly true retrieval/evidence misses. The frozen freshness audit is exhaustive because the slice has only two remaining failures: one volatile update where retrieval evidence remained stale, and one where temporary adaptation should have been used.

The strongest limitation is benchmark scope. The freshness result uses a deterministic bundled local \texttt{FreshQA}-style snapshot rather than a public benchmark run, and both experiments use one small open model. We therefore make narrow claims, release the artifact bundle, and treat broader public-benchmark and cross-model validation as future work.

\section{Artifact Availability}

The paper package is backed by frozen benchmark bundles and support artifacts listed in \texttt{paper\_run\_manifest\_index.json}. The TMLR submission uses only the official bundles under \texttt{artifacts/tmlr\_official}. The named arXiv version additionally cites one deterministic bundled controlled-update artifact, also recorded in the claims map. Dataset snapshot hashes and preprocessing provenance are summarized in \texttt{DATASET\_MANIFEST.md}. The benchmark-specific trust materials include the public FreshQA manual audit bundle together with the MQuAKE and UniEdit reliability notes.

